% Generated by the blueprint tool (blueprint-pdf/1) from a blueprint snapshot; edit the
% Markdown sources, not this file.  Compile with `latexmk -pdf` or by running
% `pdflatex` twice (the second run settles the contents and the page references).
% Images, if any, are read from assets/ next to this file.
\documentclass[11pt]{article}
\usepackage[T1]{fontenc}
\usepackage[utf8]{inputenc}
\usepackage{lmodern}
\usepackage{textcomp}
\usepackage[a4paper,margin=2.5cm,headheight=14pt]{geometry}
\usepackage{amsmath,amssymb,amsthm}
\usepackage{amscd}  % KaTeX's CD diagrams; part of amsmath, so always there
\usepackage{mathtools}
\usepackage{bm}
\usepackage{graphicx}
\usepackage{xcolor}
\usepackage{array,booktabs,tabularx}
\usepackage{fancyvrb}
\usepackage[normalem]{ulem}
\usepackage{microtype}
% Optional, and only with their fonts: a small TeX installation has the
% packages without the fonts, which fails at the first formula.
\IfFileExists{mathrsfs.sty}{\IfFileExists{rsfs10.tfm}{\usepackage{mathrsfs}}{}}{}
\IfFileExists{stmaryrd.sty}{\IfFileExists{stmary10.tfm}{\usepackage{stmaryrd}}{}}{}
\IfFileExists{cancel.sty}{\usepackage{cancel}}{}
\usepackage[hyphens]{url}
\usepackage[pdfusetitle,bookmarksnumbered=false]{hyperref}
\usepackage{bookmark}
\definecolor{bplink}{HTML}{1F4E9C}
\definecolor{bpmuted}{HTML}{5F676E}
\hypersetup{colorlinks=true,linkcolor=bplink,urlcolor=bplink,citecolor=bplink}
\setlength{\parindent}{0pt}
\setlength{\parskip}{0.55em plus 0.2em minus 0.1em}
% Lists as tight as the page's.
\makeatletter
\def\@listi{\leftmargin\leftmargini\parsep 0pt\topsep 0.3em plus 0.1em\itemsep 0.1em plus 0.1em}
\let\@listI\@listi
\def\@listii{\leftmargin\leftmarginii\labelwidth\leftmarginii\advance\labelwidth-\labelsep\parsep 0pt\topsep 0.1em\itemsep 0.1em}
% An ordered list that does not start at 1.
\newcommand{\bpstartat}[1]{\setcounter{\@enumctr}{\numexpr#1-1\relax}}
\makeatother
\setcounter{tocdepth}{3}
\emergencystretch=3em
\allowdisplaybreaks

% A numbered statement, set as a paper sets one: "Theorem 2.1 (Title)."
% #1 is the lead ("Theorem~2.1"), #2 what follows it in roman.
\newenvironment{bpstatement}[2]{\par\addvspace{0.6\baselineskip}\noindent\textbf{#1}#2\textbf{.}\enspace\ignorespaces}
  {\par\addvspace{0.4\baselineskip}}
\newenvironment{bpstatementit}[2]{\begin{bpstatement}{#1}{#2}\itshape\ignorespaces}{\end{bpstatement}}
% A step: a relation with prose of its own, read under the object it starts from.
\newenvironment{bpstep}[1]{\par\addvspace{0.3\baselineskip}\begin{list}{}{\leftmargin=1.6em\rightmargin=0pt\topsep=0pt}\item\relax
  {\small\color{bpmuted}#1}\par\nopagebreak}{\end{list}\par}
% Notes about where an object is written out.
\newenvironment{bpnote}{\par\small\color{bpmuted}\itshape}{\par}
\newcommand{\bpomitted}[1]{\par{\small\color{bpmuted}\itshape #1\par}}
\newcommand{\bpbodyheading}[1]{\par\addvspace{0.4\baselineskip}\noindent\textbf{#1}\par\nopagebreak}
\newcommand{\bprule}{\par\medskip\noindent\hfil\rule{0.3\linewidth}{0.4pt}\hfil\par\medskip}
\newcommand{\bpboundary}[1]{\par\nopagebreak{\small\color{bpmuted}#1\par}\nopagebreak}
% A reference to something this export leaves out, and one that resolves nowhere.
\newcommand{\bpout}[1]{#1\textsuperscript{\hyperlink{bp-outside}{\ensuremath{\ast}}}}
\newcommand{\bpbroken}[1]{\textcolor{red!70!black}{[#1?]}}
\newcommand{\bpslug}[1]{\texttt{#1}}
% A bibliography entry: label, anchor, entry; the label hangs in the margin.
\newcommand{\bpbibitem}[3]{\par\setbox0\hbox{[#1]\enspace}\ifdim\wd0<5em\setbox0\hbox to 5em{[#1]\hfil}\fi
  \hangindent=\wd0 \hangafter=1 \noindent\box0\phantomsection\label{#2}#3\par\smallskip}
% A character the export knows no way to set.
\newcommand{\bpglyph}[1]{\fbox{\scriptsize #1}}
% Images, at their own size unless they are wider than the line.
\newsavebox{\bpimagebox}
\newcommand{\bpimage}[1]{\sbox{\bpimagebox}{\includegraphics{#1}}%
  \ifdim\wd\bpimagebox>\linewidth\resizebox{\linewidth}{!}{\usebox{\bpimagebox}}\else\usebox{\bpimagebox}\fi}
\newcommand{\bpmissingimage}[1]{\fbox{\parbox{0.8\linewidth}{\centering\small\color{bpmuted}Image not included: #1}}}
\DefineVerbatimEnvironment{bpverbatim}{Verbatim}{fontsize=\small,xleftmargin=1em}
\renewcommand{\qedsymbol}{\ensuremath{\blacksquare}}
\hypersetup{pdftitle={Weil Conjectures — 19 Proof of the Weil conjectures},pdfsubject={Chapter 19 Proof of the Weil conjectures; all levels},pdfcreator={blueprint}}
% KaTeX commands the website knows
\providecommand{\R}{\mathbb{R}}\providecommand{\N}{\mathbb{N}}\providecommand{\Z}{\mathbb{Z}}
\providecommand{\Q}{\mathbb{Q}}\providecommand{\C}{\mathbb{C}}
\providecommand{\Reals}{\mathbb{R}}\providecommand{\reals}{\mathbb{R}}\providecommand{\Complex}{\mathbb{C}}
\providecommand{\cnums}{\mathbb{C}}\providecommand{\natnums}{\mathbb{N}}\providecommand{\Bbb}[1]{\mathbb{#1}}
\providecommand{\bold}[1]{\mathbf{#1}}\providecommand{\frak}[1]{\mathfrak{#1}}\providecommand{\mathbbm}[1]{\mathbb{#1}}
\providecommand{\lang}{\langle}\providecommand{\rang}{\rangle}\providecommand{\bra}[1]{\langle #1\rvert}
\providecommand{\ket}[1]{\lvert #1\rangle}\providecommand{\braket}[1]{\langle #1\rangle}
\providecommand{\set}[1]{\{#1\}}\providecommand{\Set}[1]{\left\{#1\right\}}
\providecommand{\infin}{\infty}\providecommand{\isin}{\in}\providecommand{\plusmn}{\pm}\providecommand{\sub}{\subset}
\providecommand{\sube}{\subseteq}\providecommand{\supe}{\supseteq}\providecommand{\exist}{\exists}
\providecommand{\bull}{\bullet}\providecommand{\larr}{\leftarrow}\providecommand{\rarr}{\rightarrow}
\providecommand{\lrarr}{\leftrightarrow}\providecommand{\Larr}{\Leftarrow}\providecommand{\Rarr}{\Rightarrow}
\providecommand{\Lrarr}{\Leftrightarrow}\providecommand{\uarr}{\uparrow}\providecommand{\darr}{\downarrow}
\providecommand{\harr}{\leftrightarrow}\providecommand{\lArr}{\Leftarrow}\providecommand{\rArr}{\Rightarrow}
\providecommand{\hArr}{\Leftrightarrow}\providecommand{\alef}{\aleph}\providecommand{\alefsym}{\aleph}
\providecommand{\real}{\Re}\providecommand{\image}{\Im}\providecommand{\weierp}{\wp}
\providecommand{\Alpha}{\mathrm{A}}\providecommand{\Beta}{\mathrm{B}}\providecommand{\Epsilon}{\mathrm{E}}
\providecommand{\Zeta}{\mathrm{Z}}\providecommand{\Eta}{\mathrm{H}}\providecommand{\Iota}{\mathrm{I}}
\providecommand{\Kappa}{\mathrm{K}}\providecommand{\Mu}{\mathrm{M}}\providecommand{\Nu}{\mathrm{N}}
\providecommand{\Omicron}{\mathrm{O}}\providecommand{\Rho}{\mathrm{P}}\providecommand{\Tau}{\mathrm{T}}
\providecommand{\Chi}{\mathrm{X}}\providecommand{\omicron}{o}\providecommand{\gt}{>}\providecommand{\lt}{<}
\providecommand{\And}{\&}\providecommand{\hdashline}{\hline}\providecommand{\operatornamewithlimits}[1]{\operatorname*{#1}}
\providecommand{\llbracket}{[\![}\providecommand{\rrbracket}{]\!]}\providecommand{\cancel}[1]{#1}
\providecommand{\mathscr}{\mathcal}
% the characters beyond ASCII this document uses
\DeclareUnicodeCharacter{2161}{\textup{II}}
\DeclareUnicodeCharacter{2162}{\textup{III}}
% the project's macros ([katex.macros] in blueprint.toml)
\providecommand{\A}{}\renewcommand{\A}{\mathbf A}
\providecommand{\AEn}{}\renewcommand{\AEn}{A_{E,n}}
\providecommand{\Aaf}{}\renewcommand{\Aaf}{\mathbf A}
\providecommand{\Ab}{}\renewcommand{\Ab}{\mathsf{Ab}}
\providecommand{\AffSch}{}\renewcommand{\AffSch}{\operatorname{AffSch}}
\providecommand{\Aone}{}\renewcommand{\Aone}{\mathbf A^1}
\providecommand{\Bl}{}\renewcommand{\Bl}{\operatorname{Bl}}
\providecommand{\Br}{}\renewcommand{\Br}{\mathrm{Br}}
\providecommand{\C}{}\renewcommand{\C}{\mathbb C}
\providecommand{\Cech}{}\renewcommand{\Cech}{\check{H}}
\providecommand{\Coext}{}\renewcommand{\Coext}{\operatorname{Coext}}
\providecommand{\Cone}{}\renewcommand{\Cone}{\operatorname{Cone}}
\providecommand{\Cons}{}\renewcommand{\Cons}{\operatorname{Cons}}
\providecommand{\Cref}{}\renewcommand{\Cref}[1]{\text{#1}}
\providecommand{\D}{}\renewcommand{\D}{D}
\providecommand{\DLAFT}{}\renewcommand{\DLAFT}{\operatorname{FT}}
\providecommand{\DLAFr}{}\renewcommand{\DLAFr}{\operatorname{Fr}}
\providecommand{\DLAQ}{}\renewcommand{\DLAQ}{{\overline{\mathbf Q}_\ell}}
\providecommand{\DLRQ}{}\renewcommand{\DLRQ}{{\overline{\mathbf Q}_\ell}}
\providecommand{\Dbc}{}\renewcommand{\Dbc}{D^b_c}
\providecommand{\Dc}{}\renewcommand{\Dc}{D_{\mathrm{cons}}}
\providecommand{\Dctf}{}\renewcommand{\Dctf}{D^b_{ctf}}
\providecommand{\Dic}{}\renewcommand{\Dic}{D_{\IndCons}}
\providecommand{\Dplus}{}\renewcommand{\Dplus}{D^{+}}
\providecommand{\End}{}\renewcommand{\End}{\operatorname{End}}
\providecommand{\Et}{}\renewcommand{\Et}{\mathrm{\acute{e}t}}
\providecommand{\Ex}{}\renewcommand{\Ex}{\mathrm{Ex}}
\providecommand{\Ext}{}\renewcommand{\Ext}{\operatorname{Ext}}
\providecommand{\F}{}\renewcommand{\F}{\mathbb F}
\providecommand{\FT}{}\renewcommand{\FT}{\operatorname{FT}}
\providecommand{\Fbar}{}\renewcommand{\Fbar}{{\overline{\mathbf F}_q}}
\providecommand{\Fell}{}\renewcommand{\Fell}{\mathbf F_\ell}
\providecommand{\Fib}{}\renewcommand{\Fib}{\operatorname{Fib}}
\providecommand{\Fix}{}\renewcommand{\Fix}{\operatorname{Fix}}
\providecommand{\Fp}{}\renewcommand{\Fp}{{\mathbf F_{p}}}
\providecommand{\Fq}{}\renewcommand{\Fq}{{\mathbf F_q}}
\providecommand{\Fqn}{}\renewcommand{\Fqn}{{\mathbf F_{q^{n}}}}
\providecommand{\Frob}{}\renewcommand{\Frob}{\operatorname{Frob}}
\providecommand{\GL}{}\renewcommand{\GL}{\operatorname{GL}}
\providecommand{\Gal}{}\renewcommand{\Gal}{\operatorname{Gal}}
\providecommand{\Hens}{}\renewcommand{\Hens}{\operatorname{Hens}}
\providecommand{\Hilb}{}\renewcommand{\Hilb}{\operatorname{Hilb}}
\providecommand{\Hom}{}\renewcommand{\Hom}{\operatorname{Hom}}
\providecommand{\IndCons}{}\renewcommand{\IndCons}{\operatorname{IndCons}}
\providecommand{\Kar}{}\renewcommand{\Kar}{\operatorname{Kar}}
\providecommand{\LT}{}\renewcommand{\LT}{\operatorname{LT}}
\providecommand{\Lie}{}\renewcommand{\Lie}{\operatorname{Lie}}
\providecommand{\Loc}{}\renewcommand{\Loc}{\operatorname{Loc}}
\providecommand{\Mod}{}\renewcommand{\Mod}{\mathsf{Mod}}
\providecommand{\N}{}\renewcommand{\N}{\mathbb{N}}
\providecommand{\Needspace}{}\renewcommand{\Needspace}{\par}
\providecommand{\O}{}\renewcommand{\O}{{\mathcal{O}}}
\providecommand{\Ob}{}\renewcommand{\Ob}{\mathrm{Ob}}
\providecommand{\Oe}{}\renewcommand{\Oe}{\mathcal O_E}
\providecommand{\PP}{}\renewcommand{\PP}{\mathbf{P}}
\providecommand{\Perf}{}\renewcommand{\Perf}{\operatorname{Perf}}
\providecommand{\Pic}{}\renewcommand{\Pic}{\mathrm{Pic}}
\providecommand{\Pm}{}\renewcommand{\Pm}{\mathbf P^{m}}
\providecommand{\Pone}{}\renewcommand{\Pone}{\mathbf P^{1}}
\providecommand{\Pqn}{}\renewcommand{\Pqn}{\mathbf P^{n}}
\providecommand{\ProFin}{}\renewcommand{\ProFin}{\operatorname{ProFin}}
\providecommand{\Proet}{}\renewcommand{\Proet}{\mathrm{pro\acute{e}t}}
\providecommand{\Q}{}\renewcommand{\Q}{\mathbb Q}
\providecommand{\Qb}{}\renewcommand{\Qb}{\overline{\mathbf Q}_{\ell}}
\providecommand{\Qbar}{}\renewcommand{\Qbar}{\overline{\mathbf Q}}
\providecommand{\QbarX}{}\renewcommand{\QbarX}{\overline{\mathbf Q}_{\ell,X}}
\providecommand{\Qbarl}{}\renewcommand{\Qbarl}{{\overline{\mathbf Q}_{\ell}}}
\providecommand{\Qell}{}\renewcommand{\Qell}{{\mathbf{Q}_{\ell}}}
\providecommand{\Ql}{}\renewcommand{\Ql}{\mathbf Q_{\ell}}
\providecommand{\Qlb}{}\renewcommand{\Qlb}{{\overline{\mathbf Q}_\ell}}
\providecommand{\Qlbar}{}\renewcommand{\Qlbar}{\overline{\mathbf Q}_\ell}
\providecommand{\R}{}\renewcommand{\R}{\mathbb{R}}
\providecommand{\RG}{}\renewcommand{\RG}{\operatorname{\mathrm{R}\Gamma}}
\providecommand{\RGamma}{}\renewcommand{\RGamma}{R\Gamma}
\providecommand{\RGammac}{}\renewcommand{\RGammac}{R\Gamma_c}
\providecommand{\RHom}{}\renewcommand{\RHom}{R\mathcal{H}om}
\providecommand{\RHomX}{}\renewcommand{\RHomX}{R\mathcal H om_X}
\providecommand{\Rat}{}\renewcommand{\Rat}{\operatorname{IsRat}}
\providecommand{\Res}{}\renewcommand{\Res}{\operatorname{Res}}
\providecommand{\Rlim}{}\renewcommand{\Rlim}{R\!\varprojlim}
\providecommand{\SIm}{}\renewcommand{\SIm}{\operatorname{SIm}}
\providecommand{\Sch}{}\renewcommand{\Sch}{\mathsf{Sch}}
\providecommand{\Sh}{}\renewcommand{\Sh}{\operatorname{Sh}}
\providecommand{\Span}{}\renewcommand{\Span}{\operatorname{span}}
\providecommand{\Spec}{}\renewcommand{\Spec}{\operatorname{Spec}}
\providecommand{\Tr}{}\renewcommand{\Tr}{\operatorname{Tr}}
\providecommand{\Trsm}{}\renewcommand{\Trsm}{\operatorname{Tr}^{\mathrm{sm}}}
\providecommand{\Wt}{}\renewcommand{\Wt}{\operatorname{Wt}}
\providecommand{\Z}{}\renewcommand{\Z}{\mathbb{Z}}
\providecommand{\Zar}{}\renewcommand{\Zar}{\mathrm{zar}}
\providecommand{\Zell}{}\renewcommand{\Zell}{{\mathbf{Z}_{\ell}}}
\providecommand{\Zl}{}\renewcommand{\Zl}{{\mathbf Z_\ell}}
\providecommand{\affet}{}\renewcommand{\affet}[1]{{#1}^{\mathrm{aff}}_{\mathrm{ét}}}
\providecommand{\affproet}{}\renewcommand{\affproet}[1]{{#1}^{\mathrm{aff}}_{\mathrm{proét}}}
\providecommand{\boxt}{}\renewcommand{\boxt}{\boxtimes_A^{\mathrm L}}
\providecommand{\cF}{}\renewcommand{\cF}{\mathcal{F}}
\providecommand{\cHom}{}\renewcommand{\cHom}{R\mathcal{H}\!om}
\providecommand{\cI}{}\renewcommand{\cI}{\mathcal{I}}
\providecommand{\cL}{}\renewcommand{\cL}{\mathcal L}
\providecommand{\cM}{}\renewcommand{\cM}{\mathcal M}
\providecommand{\cN}{}\renewcommand{\cN}{\mathcal N}
\providecommand{\cO}{}\renewcommand{\cO}{\mathcal O}
\providecommand{\cT}{}\renewcommand{\cT}{\mathcal T}
\providecommand{\checked}{}\renewcommand{\checked}{\color{red}\aftergroup\restorecolor}
\providecommand{\coeff}{}\renewcommand{\coeff}{\operatorname{coeff}}
\providecommand{\coker}{}\renewcommand{\coker}{\operatorname{coker}}
\providecommand{\colim}{}\renewcommand{\colim}{\operatorname{colim}}
\providecommand{\constructedby}{}\renewcommand{\constructedby}[1]{\par\smallskip\noindent\textcolor{purple!60!black}{\textsf{Constructed by:} }
  \footnotesize\sloppy\nolinkurl{#1}\par}
\providecommand{\deepresult}{}\renewcommand{\deepresult}[1]{\par\smallskip\noindent\textcolor{orange!45!black}{
  \textsf{Blueprint role: named deep theorem node; see the cluster proof architecture.}}
  \par\noindent\textcolor{orange!45!black}{\textsf{Mathematical source:} }#1\par}
\providecommand{\dt}{}\renewcommand{\dt}{\mathbin{\otimes^{\mathbf L}}}
\providecommand{\dw}{}\renewcommand{\dw}{\operatorname{dw}}
\providecommand{\et}{}\renewcommand{\et}[1]{{#1}_{\mathrm{ét}}}
\providecommand{\ev}{}\renewcommand{\ev}{\operatorname{ev}}
\providecommand{\fib}{}\renewcommand{\fib}{\operatorname{fib}}
\providecommand{\formalizationnote}{}\renewcommand{\formalizationnote}[1]{\begin{quote}\small\raggedright\textbf{Formalization note.} #1\end{quote}}
\providecommand{\fpqc}{}\renewcommand{\fpqc}{X_{\mathrm{fpqc}}}
\providecommand{\gammacont}{}\renewcommand{\gammacont}{\Gamma_{\mathrm{cont}}}
\providecommand{\gm}{}\renewcommand{\gm}{\mathbf G_m}
\providecommand{\hocolim}{}\renewcommand{\hocolim}{\operatorname{hocolim}}
\providecommand{\holim}{}\renewcommand{\holim}{\operatorname{holim}}
\providecommand{\hrB}{}\renewcommand{\hrB}{\@ifnextchar]{\hrEnd}{\hrStep}}
\providecommand{\hrEnd}{}\renewcommand{\hrEnd}{\@firstoftwo{}}
\providecommand{\hrStep}{}\renewcommand{\hrStep}{\@firstoftwo{\hrB}}
\providecommand{\id}{}\renewcommand{\id}{\operatorname{id}}
\providecommand{\kunFl}{}\renewcommand{\kunFl}{\mathbf F_\ell}
\providecommand{\kunZn}{}\renewcommand{\kunZn}{\mathbf Z/n\mathbf Z}
\providecommand{\kuncd}{}\renewcommand{\kuncd}{\operatorname{cd}}
\providecommand{\kundim}{}\renewcommand{\kundim}{\dim}
\providecommand{\kunotimes}{}\renewcommand{\kunotimes}[1]{\otimes^{\mathbf L}_{#1}}
\providecommand{\kunsupp}{}\renewcommand{\kunsupp}{\operatorname{Supp}}
\providecommand{\lotimes}{}\renewcommand{\lotimes}{\otimes^{\mathrm L}}
\providecommand{\m}{}\renewcommand{\m}{{\mathfrak{m}}}
\providecommand{\n}{}\renewcommand{\n}{{\mathfrak{n}}}
\providecommand{\nagataGap}{}\renewcommand{\nagataGap}[1]{\par\noindent\textbf{Formalization gap.} #1\par}
\providecommand{\openobligation}{}\renewcommand{\openobligation}[1]{\par\smallskip\noindent\textcolor{red!65!black}{\textsf{Open obligation:} }
  \footnotesize\sloppy\nolinkurl{#1}\par}
\providecommand{\p}{}\renewcommand{\p}{{\mathfrak{p}}}
\providecommand{\pr}{}\renewcommand{\pr}{\operatorname{pr}}
\providecommand{\preshv}{}\renewcommand{\preshv}[1]{\mathrm{PreShv}(#1)}
\providecommand{\pro}{}\renewcommand{\pro}[1]{\mathrm{Pro}(#1)}
\providecommand{\proet}{}\renewcommand{\proet}{X_{\mathrm{proét}}}
\providecommand{\q}{}\renewcommand{\q}{{\mathfrak{q}}}
\providecommand{\rad}{}\renewcommand{\rad}{\operatorname{rad}}
\providecommand{\rank}{}\renewcommand{\rank}{\operatorname{rank}}
\providecommand{\red}{}\renewcommand{\red}{\color{red}}
\providecommand{\restorecolor}{}\renewcommand{\restorecolor}{\color{black}}
\providecommand{\retiredobligation}{}\renewcommand{\retiredobligation}[1]{\par\smallskip\noindent\textcolor{red!75!black}{\textsf{Retired invalid route:} }
  \footnotesize\sloppy\nolinkurl{#1}\par}
\providecommand{\rk}{}\renewcommand{\rk}{\operatorname{rk}}
\providecommand{\sep}{}\renewcommand{\sep}{\textnormal{sep}}
\providecommand{\shv}{}\renewcommand{\shv}[1]{\mathrm{Shv}(#1)}
\providecommand{\spec}{}\renewcommand{\spec}[1]{\mathrm{Spec}(#1)}
\providecommand{\stacks}{}\renewcommand{\stacks}[1]{\href{https://stacks.math.columbia.edu/tag/#1}{\textnormal{Stacks, Tag~#1}}}
\providecommand{\stacksproject}{}\renewcommand{\stacksproject}[1]{{\cite[\href{https://stacks.math.columbia.edu/tag/#1}{Tag #1}]{stacks-project}}}
\providecommand{\statusnote}{}\renewcommand{\statusnote}[1]{\begin{quote}\small\raggedright\textbf{Status.} #1\end{quote}}
\providecommand{\tildeideal}{}\renewcommand{\tildeideal}[2]{\widetilde{(#1, #2)}}
\providecommand{\tildering}{}\renewcommand{\tildering}[3]{{#1}_{\widetilde{(#2, #3)}}}
\providecommand{\tr}{}\renewcommand{\tr}{\operatorname{tr}}
\providecommand{\wt}{}\renewcommand{\wt}{\operatorname{wt}}

\begin{document}
\pagestyle{plain}
\begin{center}
{\LARGE\bfseries Weil Conjectures\par}
\vspace{0.7em}
{\large Chapter 19: Proof of the Weil conjectures\par}
\vspace{0.5em}
{\small\color{bpmuted}Snapshot 2026-09-26, revision 8e4cb88e2a\par}
\end{center}
\vspace{0.5em}
{\small\tableofcontents}
\bigskip\hrule\bigskip

% object: chap-Proof-of-the-Weil-conjectures (blueprint/chap-Proof-of-the-Weil-conjectures/_section.md)
\section*{\texorpdfstring{19\quad Proof of the Weil conjectures}{19 Proof of the Weil conjectures}}
\phantomsection\label{bp1}%
\addcontentsline{toc}{section}{\texorpdfstring{19\quad Proof of the Weil conjectures}{19 Proof of the Weil conjectures}}
% text: prose-19-01
\phantomsection\label{bp2}%
The Deligne–Laumon theorem is an upper bound for weights, stated with \(\ell\)-adic coefficients and one embedding \(\iota:\overline{\mathbf Q}_\ell\to\mathbf C\). Weil's conjectures ask for more: an equality \(|\alpha|=q^{i/2}\), for polynomials \(P_i\) with integer coefficients, under every complex embedding. This chapter closes both gaps.

The lower bound comes from Poincaré duality, and a curve shows how. Let \(X_0\) be a smooth projective curve and \(\alpha\) an eigenvalue of Frobenius on \(H^1(X,\overline{\mathbf Q}_\ell)\). Then Theorem \bpout{13.11}, applied to the constant sheaf, gives \(|\iota(\alpha)|\le q^{1/2}\). On the other hand, cup product is a perfect pairing \(H^1\times H^1\to H^2=\overline{\mathbf Q}_\ell(-1)\), which Frobenius multiplies by \(q\). Hence \(q/\alpha\) is an eigenvalue as well, and the same bound for \(q/\alpha\) says that \(|\iota(\alpha)|\ge q^{1/2}\). In dimension \(d\) the argument is the same: an eigenvalue \(\alpha\) on \(H^i(X,\overline{\mathbf Q}_\ell)\) satisfies \(|\iota(\alpha)|\le q^{i/2}\), while duality makes \(q^d/\alpha\) an eigenvalue on \(H^{2d-i}\), of absolute value at most \(q^{(2d-i)/2}\), and we find \(|\iota(\alpha)|=q^{i/2}\). This equality is Weil's axiom (d), and we prove it in Section \hyperref[bp3]{19.1}, more generally for the higher direct images of a pure lisse sheaf under a smooth proper morphism (Theorem \hyperref[bp5]{19.1}); the case of the constant sheaf over a point is Corollary \hyperref[bp8]{19.2}. Since the Deligne–Laumon theorem holds for every \(\iota\), purity holds for all complex embeddings at once. Indeed, every embedding of a number field into \(\mathbf C\) extends to an isomorphism \(\overline{\mathbf Q}_\ell\simeq\mathbf C\).

The polynomials \(P_{i,\ell}=\det(1-TF\mid H^i(X,\mathbf Q_\ell))\) are defined with \(\ell\)-adic coefficients, and a priori they could depend on \(\ell\). The zeta function does not mention \(\ell\). It has a unique expression \(Z(X_0,T)=A/B\) with \(A,B\in\mathbf Z[T]\) coprime and \(A(0)=B(0)=1\) (Lemma \hyperref[bp15]{19.4}; Sections \hyperref[bp10]{19.2}–\hyperref[bp18]{Section 19.3}). Can each \(P_{i,\ell}\) be read off from \(Z(X_0,T)\)? The trace formula writes \(Z\) as the alternating product of the \(P_{i,\ell}\), but an alternating product can cancel. For example, \(X_0=\mathbf G_m\sqcup
\operatorname{Spec}\mathbf F_q\) is not proper, and it has \(q^m\) points over \(\mathbf F_{q^m}\), as many as \(\mathbf A^1\). Hence \(Z(X_0,T)=1/(1-qT)\). Yet \(H^0_c\) and \(H^1_c\) are both one-dimensional with eigenvalue \(1\), and the factor \(1-T\) cancels between them. For a smooth proper \(X_0\) purity rules this out. The reciprocal roots of \(P_{i,\ell}\) have absolute value \(q^{i/2}\), so no root of \(P_{i,\ell}\) is a root of \(P_{j,\ell}\) for \(i\neq j\). It follows that the product of the \(P_{i,\ell}\) with \(i\) odd is \(A\) and the product with \(i\) even is \(B\). Grouping the reciprocal roots of \(A\) and \(B\) by absolute value then recovers each \(P_{i,\ell}\) from \(Z(X_0,T)\) alone (Theorem \hyperref[bp24]{19.7}). In particular \(P_{i,\ell}\) has integer coefficients and does not depend on \(\ell\).

In terms of point counts, purity gives the classical estimates. For a smooth proper geometrically connected \(X_0\) we obtain \(|\#X_0(\mathbf F_{q^m})-(1+q^{dm})|\le\sum_{i=1}^{2d-1}b_iq^{mi/2}\) (Corollary \hyperref[bp30]{19.10}). In particular, for a curve of genus \(g\) this is Weil's bound \(2gq^{m/2}\), and for elliptic curves it is Hasse's theorem. As an illustration, Example \hyperref[bp28]{19.9} works this out for genus-one curves over \(\mathbf F_3\).

% object: sec-The-Riemann-hypothesis (blueprint/chap-Proof-of-the-Weil-conjectures/sec-The-Riemann-hypothesis/_section.md), content hidden
\subsection*{\texorpdfstring{19.1\quad The Riemann hypothesis}{19.1 The Riemann hypothesis}}
\phantomsection\label{bp3}%
\addcontentsline{toc}{subsection}{\texorpdfstring{19.1\quad The Riemann hypothesis}{19.1 The Riemann hypothesis}}
% text: prose-19-02
\phantomsection\label{bp4}%
For a proper map we have \(Rf_!=Rf_*\), so the upper bound of Theorem \bpout{18.15} applies to ordinary cohomology. If the map is also smooth, Poincaré duality on each fibre pairs \(H^i(X_{\bar y},F)\) perfectly with \(H^{2d-i}(X_{\bar y},F^\vee(d))\). The eigenvalues of the second group are the inverses of those of the first. The second group has weights at most \((-w-2d)+(2d-i)=-w-i\), and this upper bound becomes the lower bound \(w+i\) for the first.

Both hypotheses are needed. Without properness, \(H^1_c\) of \(\mathbf G_m\) has weight \(0\) instead of \(1\). Without smoothness duality fails. For example, the nodal cubic is proper, and its \(H^1\) is a line on which Frobenius acts by \(\pm1\), of weight \(0\). The relative statement is no harder than the absolute one, because we check the bound at one closed point \(y\) at a time, on the smooth proper fibre over \(\kappa(y)\). In the proof of Theorem \hyperref[bp24]{19.7} we use only the case of a point as base and the constant sheaf (Corollary \hyperref[bp8]{19.2}).

% object: theorem-purity-higher-direct-images-under (blueprint/chap-Proof-of-the-Weil-conjectures/sec-The-Riemann-hypothesis/theorem-purity-higher-direct-images-under/_section.md)
\phantomsection\label{bp5}%
\begin{bpstatementit}{Theorem~19.1}{ (Purity of higher direct images under a smooth proper morphism)}
Let \(f_0:X_0\to Y_0\) be smooth proper of pure relative dimension \(d\), and let \(F\) be a lisse Weil sheaf, punctually \(\iota\)-pure of weight \(w\). For every \(i\), the canonical identification

\[
R^if_{0!}F=R^if_{0*}F
\]

is Weil-equivariant, and this constructible sheaf is punctually \(\iota\)-pure of weight \(w+i\). In particular, it is mixed of the single weight \(w+i\).
\end{bpstatementit}

% proof: proof/theorem-purity-higher-direct-images-under
\begin{proof}[{Proof}]\phantomsection\label{bp6}%
Since \(f_0\) is proper, the definition of \(Rf_{0!}\) by compactification identifies it with \(Rf_{0*}\). This identification commutes with the base change maps that define the Weil structures, which proves the first assertion.

Fix a closed point \(y\). By proper base change, the stalk at \(\bar y\) with its local Frobenius is \(H^i(X_{\bar y},F)\). By Theorem \bpout{18.15}, \(R^if_{0*}F\) is mixed of weights at most \(w+i\), so every eigenvalue \(\alpha\) there satisfies

\[
|\iota(\alpha)|\leq q_y^{(w+i)/2}.
\]

The opposite inequality comes from duality. The dual \(F^\vee\) is lisse of weight \(-w\), so \(F^\vee(d)\) has weight \(-w-2d\). On the smooth proper geometric fibre, Corollary \bpout{10.13} and Theorem \bpout{10.21} give a Frobenius-equivariant perfect pairing

\[
H^i(X_{\bar y},F)\times
H^{2d-i}(X_{\bar y},F^\vee(d))\longrightarrow\overline{\mathbf Q}_\ell.
\]

Since the pairing is perfect and takes values in the trivial Weil line, it identifies the second factor with the dual of the first as a Frobenius module, so the eigenvalues on the second factor are the \(\alpha^{-1}\), with multiplicities. Applying Theorem \bpout{18.15} once more, now to the structure map of the fibre over \(\kappa(y)\) and to \(F^\vee(d)\), of weight \(-w-2d\), and to its cohomology in degree \(2d-i\), we get

\[
|\iota(\alpha)|^{-1}
\leq q_y^{(-w-2d+2d-i)/2}=q_y^{(-w-i)/2}.
\]

The two inequalities give \(|\iota(\alpha)|=q_y^{(w+i)/2}\) at every closed point \(y\). Every nonzero pure quotient in a mixed filtration of \(R^if_{0*}F\) has a nonzero stalk at some closed point, so its weight is \(w+i\). This proves the last assertion.
\end{proof}

% text: prose-19-03
\phantomsection\label{bp7}%
For the Weil conjectures we need only the case where the base is a point and the sheaf is constant. Then the fibre is \(X\) itself, and the stalk of \(R^if_*\overline{\mathbf Q}_\ell\) is \(H^i(X,\overline{\mathbf Q}_\ell)\).

% object: corollary-purity-cohomology-smooth-proper-variety (blueprint/chap-Proof-of-the-Weil-conjectures/sec-The-Riemann-hypothesis/corollary-purity-cohomology-smooth-proper-variety/_section.md)
\phantomsection\label{bp8}%
\begin{bpstatementit}{Corollary~19.2}{ (Purity of the cohomology of a smooth proper variety)}
Let \(X_0/\mathbf F_q\) be smooth proper of pure dimension \(d\), put \(X=X_0\times_{\mathbf F_q}\overline{\mathbf F}_q\), let \(\ell\nmid q\) be a prime, and let \(\iota:\overline{\mathbf Q}_\ell\hookrightarrow\mathbf C\) be a field embedding. Then every geometric Frobenius eigenvalue \(\alpha\) on \(H^i(X_{\mathrm{pro}\acute{e}t},\overline{\mathbf Q}_\ell)\) satisfies

\[
|\iota(\alpha)|=q^{i/2}.
\]
\end{bpstatementit}

% proof: proof/corollary-purity-cohomology-smooth-proper-variety
\begin{proof}[{Proof}]\phantomsection\label{bp9}%
The constant sheaf has weight zero, since its local Frobenius is the identity. So Theorem \hyperref[bp5]{19.1} applies to the structure map \(f_0:X_0\to\operatorname{Spec}\mathbf F_q\) with \(w=0\). The base has a single closed point, with \(q_y=q\). The stalk of \(R^if_{0*}\overline{\mathbf Q}_\ell\) there is \(H^i(X,\overline{\mathbf Q}_\ell)\) with its geometric Frobenius, and its eigenvalues therefore satisfy \(|\iota(\alpha)|=q^{i/2}\).
\end{proof}


% object: sec-Rational-power-series (blueprint/chap-Proof-of-the-Weil-conjectures/sec-Rational-power-series/_section.md), content hidden
\subsection*{\texorpdfstring{19.2\quad Rational power series}{19.2 Rational power series}}
\phantomsection\label{bp10}%
\addcontentsline{toc}{subsection}{\texorpdfstring{19.2\quad Rational power series}{19.2 Rational power series}}
% text: prose-19-04
\phantomsection\label{bp11}%
The zeta function has integer coefficients by its definition, but the trace formula proves its rationality only over \(\mathbf Q_\ell\), as an alternating product of \(\ell\)-adic characteristic polynomials. To compare these polynomials for different primes \(\ell\), we need a form of the zeta function that does not mention \(\ell\). We show that a rational power series with integer coefficients is a quotient of integral polynomials in a unique reduced form.

Lemma \hyperref[bp12]{19.3} shows that rationality is a condition of linear algebra on the coefficients, namely that the Hankel matrix \((c_{r+s})\) has finite rank. Such a condition descends from \(\mathbf Q_\ell\) to \(\mathbf Q\). For integrality of the reduced fraction, Lemma \hyperref[bp15]{19.4} uses a lattice in the row span of the Hankel matrix that is stable under the shift. The characteristic polynomial of the shift on this lattice is then integral. For example, the coefficients \(c_n=2^{n+1}-1\) of \(1/((1-2T)(1-T))\) satisfy \(c_{n+2}=3c_{n+1}-2c_n\). The shift matrix \(S\) on the row span has \(\det(1-TS)=(1-2T)(1-T)\), which is the reduced denominator. In the next section we compare the \(\ell\)-adic factorizations with the reduced fraction \(A/B\).

% object: lemma-hankel-rank-rationality (blueprint/chap-Proof-of-the-Weil-conjectures/sec-Rational-power-series/lemma-hankel-rank-rationality/_section.md)
\phantomsection\label{bp12}%
\begin{bpstatementit}{Lemma~19.3}{ (Hankel rank and rationality)}
For a field \(K\) and \(C(T)=\sum_{n\geq0}c_nT^n\in K[[T]]\), the following conditions are equivalent:

\begin{enumerate}
\item \(C(T)\) is rational and regular at zero;
\item the finite submatrices of \((c_{r+s})_{r,s\geq0}\) have bounded rank;
\item the span of the rows \(R_j=(c_{j+s})_{s\geq0}\) is finite-dimensional.
\end{enumerate}

If \(K\subset K'\) and \(C\) has coefficients in \(K\), rationality over \(K'\) implies rationality over \(K\).
\end{bpstatementit}

% proof: proof/lemma-hankel-rank-rationality
\begin{proof}[{Proof}]\phantomsection\label{bp13}%
The equivalence of (2) and (3) is linear algebra. For the equivalence of (1) and (3), note that a denominator \(Q\) gives a linear recurrence among the rows. Conversely, the characteristic polynomial of the shift on a finite-dimensional row span is a denominator. The descent statement holds because (2) does not change when we extend the field.

For example, \(C=1/(1-T)\) has all \(c_n=1\), every row is \((1,1,\ldots)\), and the row span is a line. For \(C=1/(1-T)^2\) we have \(c_n=n+1\), and \(R_j=R_0+j(R_1-R_0)\), so the row span is a plane.

For (3)\(\Rightarrow\)(2), a nonzero \(k\times k\) minor of a finite submatrix gives \(k\) linearly independent rows. So every finite submatrix has rank at most the dimension of the row span. For (2)\(\Rightarrow\)(3), let \(R_{j_1},\ldots,R_{j_m}\) be linearly independent. On their span the coordinate functionals have common kernel zero. They therefore span its dual, and \(m\) of them already do. The corresponding \(m\times m\) minor is nonzero, and so \(m\) is at most the bound in (2).

For (1)\(\Rightarrow\)(3), suppose that \(QC=P\) with \(Q(0)=1\). Comparing coefficients gives a fixed linear recurrence for all sufficiently large indices, and the same recurrence expresses every sufficiently late row in terms of a fixed finite list of earlier rows.

For (3)\(\Rightarrow\)(1), the left shift \(S(v)_n=v_{n+1}\) preserves the finite-dimensional row span \(V\), since \(S(R_j)=R_{j+1}\). Let \(e\) be evaluation at coordinate zero and \(v=R_0\). In any basis of \(V\), with \(M\) the matrix of the shift, we have

\[
C(T)=e(I-TM)^{-1}v
=\frac{e\operatorname{adj}(I-TM)v}{\det(I-TM)}.
\]

The first equality holds coefficient by coefficient. Indeed, \((I-TM)^{-1}=\sum_nT^nM^n\), and \(eM^nv\) is the coordinate zero of \(S^n(R_0)=R_n\), that is, \(c_n\). As the denominator has constant coefficient one, this proves (1).

For the descent, note that condition (2) involves only the coefficients \(c_n\in K\), and the rank of a matrix with entries in \(K\) does not change when we pass to \(K'\). So (2) holds over \(K\) if and only if it holds over \(K'\). Applying (1)\(\Rightarrow\)(2) over \(K'\) and (2)\(\Rightarrow\)(1) over \(K\) gives the descent of rationality.
\end{proof}

% text: prose-19-05
\phantomsection\label{bp14}%
Applied with \(K=\mathbf Q\) and \(K'=\mathbf Q_\ell\), the lemma shows that the zeta function, which the trace formula writes as a quotient of polynomials over \(\mathbf Q_\ell\), is a quotient of polynomials over \(\mathbf Q\). The next lemma shows that the integrality of its coefficients passes to the reduced fraction.

% object: lemma-reduced-fractions-integer-coefficients (blueprint/chap-Proof-of-the-Weil-conjectures/sec-Rational-power-series/lemma-reduced-fractions-integer-coefficients/_section.md)
\phantomsection\label{bp15}%
\begin{bpstatementit}{Lemma~19.4}{ (Reduced fractions with integer coefficients)}
If \(C(T)\in1+T\mathbf Z[[T]]\) is rational over \(\mathbf Q\), then there are unique coprime \(A,B\in\mathbf Z[T]\) such that

\[
C(T)=A(T)/B(T),\qquad A(0)=B(0)=1.
\]
\end{bpstatementit}

% proof: proof/lemma-reduced-fractions-integer-coefficients
\begin{proof}[{Proof}]\phantomsection\label{bp16}%
We repeat the construction in the proof of Lemma \hyperref[bp12]{19.3} with an integral basis of the row span. This gives an integral fraction, which we then reduce with Gauss' lemma.

Let \(V\) be the rational row span of Lemma \hyperref[bp12]{19.3}, and put \(L=V\cap\mathbf Z^{\mathbf Z_{\geq0}}\). Then \(L\) contains all rows, spans \(V\), and is preserved by the left shift. Since \(V\) is finite-dimensional, some \(r\) of the coordinate functionals span its dual, as in the proof of (2)\(\Rightarrow\)(3) in Lemma \hyperref[bp12]{19.3}, and so they are jointly injective on \(V\). These coordinates embed \(L\) into \(\mathbf Z^r\). Since a subgroup of \(\mathbf Z^r\) is free of finite rank, \(L\) is a shift-stable lattice in \(V\), free of rank \(\dim V\).

In a basis of \(L\), the shift matrix \(M\), the row \(e\) and the column \(v=R_0\) are integral. The matrix identity in the proof of Lemma \hyperref[bp12]{19.3} yields integral polynomials

\[
A_0=e\operatorname{adj}(I-TM)v,\qquad B_0=\det(I-TM),
\qquad C=A_0/B_0.
\]

Both constant coefficients are one. Let \(G\) be a primitive integral polynomial that is a greatest common divisor of \(A_0\) and \(B_0\) in \(\mathbf Q[T]\). By Gauss' lemma the quotients \(A_0/G\) and \(B_0/G\) are integral. Since \(G(0)\) divides \(B_0(0)=1\), we may choose the sign of \(G\) so that \(G(0)=1\). Dividing by \(G\), we obtain polynomials \(A\) and \(B\) that are coprime and have constant coefficients one.

For uniqueness, let \(A/B=A'/B'\) be two such fractions, so that \(AB'=A'B\) in \(\mathbf Q[T]\). Since \(B\) is coprime to \(A\), it divides \(B'\), and in the same way \(B'\) divides \(B\). So \(B'=cB\) for a scalar \(c\). Comparing constant coefficients gives \(c=1\), and then \(AB=A'B\) forces \(A=A'\).
\end{proof}

% text: prose-19-06
\phantomsection\label{bp17}%
By Corollary \bpout{12.31}, \(Z(X_0,T)\in\mathbf Q(T)\cap(1+T\mathbf Z[[T]])\) for every separated \(X_0\) of finite type over \(\mathbf F_q\). So Lemma \hyperref[bp15]{19.4} applies. We write \(Z(X_0,T)=A/B\) for the resulting fraction, with \(A\) and \(B\) coprime in \(\mathbf Z[T]\) and \(A(0)=B(0)=1\).


% object: sec-Integrality-and-independence-of-ell (blueprint/chap-Proof-of-the-Weil-conjectures/sec-Integrality-and-independence-of-ell/_section.md), content hidden
\subsection*{\texorpdfstring{19.3\quad Integrality and independence of $\ell$}{19.3 Integrality and independence of ell}}
\phantomsection\label{bp18}%
\addcontentsline{toc}{subsection}{\texorpdfstring{19.3\quad Integrality and independence of $\ell$}{19.3 Integrality and independence of ell}}
% text: prose-19-07
\phantomsection\label{bp19}%
Let \(X_0\) be smooth and proper over \(\mathbf F_q\), of pure dimension \(d\), and let \(A/B\) be the reduced fraction of \(Z(X_0,T)\) (Lemma \hyperref[bp15]{19.4}). We now carry out the argument sketched at the start of the chapter. The point is that purity makes the characteristic polynomials of Frobenius in different degrees pairwise coprime, so no cancellation can occur, and the polynomials can be read off from \(A\) and \(B\). For a prime \(\ell\nmid q\) put

\[
H^i_\ell=H^i(X_{\mathrm{pro}\acute{e}t},\mathbf Q_\ell),\qquad
P_{i,\ell}(T)=\det(1-TF\mid H^i_\ell),\qquad
b_{i,\ell}=\dim_{\mathbf Q_\ell}H^i_\ell.
\]

Since \(X_0\) is proper, \(H^i_c=H^i\), and the determinant identity of Corollary \bpout{12.31} for \(\Lambda=\mathbf Q_\ell\) reads

\[
Z(X_0,T)=\prod_{i=0}^{2d}P_{i,\ell}(T)^{(-1)^{i+1}}.
\tag{19.1}
\]

Since Frobenius is invertible, \(\deg P_{i,\ell}=b_{i,\ell}\).

% object: definition-reciprocal-root (blueprint/chap-Proof-of-the-Weil-conjectures/sec-Integrality-and-independence-of-ell/definition-reciprocal-root/_section.md)
\phantomsection\label{bp20}%
\begin{bpstatement}{Definition~19.5}{ (Reciprocal roots)}
A \emph{reciprocal root} of a polynomial \(P\) over a field is an element \(\alpha\) of an algebraic closure such that \(1-\alpha T\) divides \(P\).
\end{bpstatement}

% text: prose-19-08
\phantomsection\label{bp21}%
Thus the reciprocal roots of \(P_{i,\ell}\) are the Frobenius eigenvalues on \(H^i_\ell\), counted with multiplicity. The next lemma shows that \(A\) and \(B\) sort these eigenvalues by degree, independently of \(\ell\) and of the complex embedding.

% object: lemma-reciprocal-roots-zeta (blueprint/chap-Proof-of-the-Weil-conjectures/sec-Integrality-and-independence-of-ell/lemma-reciprocal-roots-zeta/_section.md)
\phantomsection\label{bp22}%
\begin{bpstatementit}{Lemma~19.6}{ (Reciprocal roots of the zeta function)}
Let \(X_0/\mathbf F_q\) be smooth proper of pure dimension \(d\), let \(Z(X_0,T)=A/B\) be the reduced fraction fixed above, and let \(\ell\nmid q\) be a prime. Then

\[
A=\prod_{i\ \mathrm{odd}}P_{i,\ell},\qquad
B=\prod_{i\ \mathrm{even}}P_{i,\ell}
\quad\hbox{in }\mathbf Q_\ell[T].
\tag{19.2}
\]

Every reciprocal root \(a\in\overline{\mathbf Q}\) of \(A\) or of \(B\) is a nonzero algebraic integer, and there is an integer \(i\), odd for \(A\) and even for \(B\), such that \(|\tau(a)|=q^{i/2}\) for every field embedding \(\tau:\mathbf Q(a)\hookrightarrow\mathbf C\) and such that \(\theta(a)\) is a Frobenius eigenvalue on \(H^i_\ell\) for every field embedding \(\theta:\overline{\mathbf Q}\hookrightarrow\overline{\mathbf Q}_\ell\).
\end{bpstatementit}

% proof: proof/lemma-reciprocal-roots-zeta
\begin{proof}[{Proof}]\phantomsection\label{bp23}%
The identity (19.2) holds because both sides are reduced fractions for \(Z(X_0,T)\) over \(\mathbf Q_\ell\). By Corollary \hyperref[bp8]{19.2}, the reciprocal roots of \(P_{i,\ell}\) have absolute value \(q^{i/2}\) under every embedding \(\iota:\overline{\mathbf Q}_\ell\hookrightarrow\mathbf C\). Since \(q>1\), distinct degrees give distinct absolute values, and so the polynomials \(P_{i,\ell}\) are pairwise coprime. Their products over odd and over even degrees are therefore coprime, and they have constant coefficients one. By (19.1) their quotient is \(Z(X_0,T)\). On the other hand, \(A\) and \(B\) remain coprime in \(\mathbf Q_\ell[T]\), since a Bézout relation over \(\mathbf Q\) remains valid over \(\mathbf Q_\ell\). The two reduced fractions agree by the uniqueness argument in the proof of Lemma \hyperref[bp15]{19.4}, which works over any field.

Write \(A(T)=\prod_a(1-aT)\) and \(B(T)=\prod_b(1-bT)\) over \(\overline{\mathbf Q}\). The reciprocal roots \(a\) and \(b\) are roots of the monic integral polynomials \(U^{\deg A}A(U^{-1})\) and \(U^{\deg B}B(U^{-1})\). Hence they are nonzero algebraic integers. Let \(\theta:\overline{\mathbf Q}\hookrightarrow\overline{\mathbf Q}_\ell\) be an embedding, and let \(a\) be a reciprocal root of \(A\) or \(B\). By (19.2) and the pairwise coprimality of the \(P_{i,\ell}\), there is exactly one \(i\) such that \(\theta(a)\) is a Frobenius eigenvalue on \(H^i_\ell\). It is odd for \(A\) and even for \(B\).

Corollary \hyperref[bp8]{19.2} bounds \(|\iota(\theta(a))|\) for embeddings \(\iota\) of \(\overline{\mathbf Q}_\ell\). The lemma, however, concerns the embeddings of the number field \(\mathbf Q(a)\). We pass from one to the other by extending field embeddings. Let \(\tau:\mathbf Q(a)\hookrightarrow\mathbf C\) be any embedding. The map \(\tau\circ\theta^{-1}\), defined on \(\theta(\mathbf Q(a))\), extends to a field isomorphism \(\iota':\overline{\mathbf Q}_\ell\xrightarrow{\sim}\mathbf C\), because both are algebraically closed fields of characteristic zero with transcendence degree \(2^{\aleph_0}\) over \(\mathbf Q\). We have \(\iota'(\theta(a))=\tau(a)\). Applying Corollary \hyperref[bp8]{19.2} with this \(\iota'\) to the eigenvalue \(\theta(a)\) on \(H^i_\ell\), we find

\[
|\tau(a)|=q^{i/2}.
\]

Since \(q>1\), this absolute value determines \(i\), and so \(i\) does not depend on \(\theta\).
\end{proof}

% object: theorem-weil-conjectures-smooth-proper-schemes (blueprint/chap-Proof-of-the-Weil-conjectures/sec-Integrality-and-independence-of-ell/theorem-weil-conjectures-smooth-proper-schemes/_section.md)
\phantomsection\label{bp24}%
\begin{bpstatementit}{Theorem~19.7}{ (The Weil conjectures for smooth proper schemes)}
Let \(X_0/\mathbf F_q\) be smooth proper of pure dimension \(d\), and let \(P_{i,\ell}\) be as above.

\begin{enumerate}
\item There are unique polynomials \(P_i(T)\in\mathbf Z[T]\) with \(P_i(0)=1\) whose images in \(\mathbf Q_\ell[T]\) are \(P_{i,\ell}(T)\) for every prime \(\ell\nmid q\), and \(Z(X_0,T)=\prod_{i=0}^{2d}P_i(T)^{(-1)^{i+1}}\).
\item The reciprocal roots of \(P_i\) are algebraic integers, and each of them has absolute value \(q^{i/2}\) under every complex embedding.
\item The degrees \(b_i=\deg P_i=\dim_{\mathbf Q_\ell}H^i_\ell\) do not depend on \(\ell\), and \(b_i=b_{2d-i}\). If \(X_0\) is geometrically connected, then \(P_0=1-T\), \(P_{2d}=1-q^dT\), and \(b_0=b_{2d}=1\).
\item Put \(\chi=\sum_i(-1)^ib_i\). Then \(d\chi\) is even and there is \(\varepsilon\in\{1,-1\}\) such that \[
   Z\!\left(X_0,\frac1{q^dT}\right)
   =\varepsilon q^{d\chi/2}T^\chi Z(X_0,T)
   \quad\hbox{in }\mathbf Q(T).
   \]
\end{enumerate}
\end{bpstatementit}

% proof: proof/theorem-weil-conjectures-smooth-proper-schemes
\begin{proof}[{Proof}]\phantomsection\label{bp25}%
For (1) and (2), we define \(P_i\) using only the reduced fraction \(A/B\) of Lemma \hyperref[bp15]{19.4}, by sorting the reciprocal roots of \(A\) and \(B\) by absolute value, and then we compare \(P_i\) with \(P_{i,\ell}\) by Lemma \hyperref[bp22]{19.6}.

For odd \(i\), let \(R_i\) be the submultiset of the reciprocal roots of \(A\) whose absolute value under every complex embedding is \(q^{i/2}\), and for even \(i\) define \(R_i\) in the same way using the reciprocal roots of \(B\). By Lemma \hyperref[bp22]{19.6}, every reciprocal root of \(A\) or \(B\) lies in exactly one \(R_i\). The multisets of reciprocal roots of \(A\) and \(B\) are Galois-stable, with multiplicities. The condition on all complex embeddings is Galois-invariant, and so each \(R_i\) is Galois-stable. Put

\[
P_i(T)=\prod_{a\in R_i}(1-aT).
\]

The coefficients of \(P_i\) are integers, being rational by Galois invariance and algebraic integers because the roots satisfy monic integral equations.

Now let \(\ell\nmid q\) be any prime and \(\theta:\overline{\mathbf Q}\hookrightarrow\overline{\mathbf Q}_\ell\) an embedding. By (19.2) and Lemma \hyperref[bp22]{19.6}, the embedding \(\theta\) maps \(R_i\) onto the multiset of Frobenius eigenvalues on \(H^i_\ell\). Since both polynomials have constant coefficient one, \(P_i=P_{i,\ell}\) in \(\overline{\mathbf Q}_\ell[T]\), hence in \(\mathbf Q_\ell[T]\). As \(\ell\) was arbitrary, this proves the existence in (1), and uniqueness follows by comparing with a single \(P_{i,\ell}\). The product formula for \(Z(X_0,T)\) in (1) is then (19.1), and assertion (2) is Lemma \hyperref[bp22]{19.6}.

The remaining assertions follow from Poincaré duality and the functional equation of Chapter \bpout{12}. For (3), since \(P_i=P_{i,\ell}\), we have \(b_i=\deg P_i=\dim_{\mathbf Q_\ell}H^i_\ell\), which does not depend on \(\ell\), and \(b_i=b_{2d-i}\) by Corollary \bpout{12.33}.

For the second part of (3), assume that \(X_0\) is geometrically connected, so that \(X\) is connected. By Theorem \bpout{10.15}, \(H^0(X,\mathbf Q_\ell)=\mathbf Q_\ell\), with Frobenius acting as the identity, and the trace map \(\operatorname{tr}_X\) identifies \(H^{2d}(X,\mathbf Q_\ell)\) with \(\mathbf Q_\ell(-d)\), on which Frobenius acts by \(q^d\). So \(b_0=b_{2d}=1\), \(P_0=1-T\) and \(P_{2d}=1-q^dT\).

For (4), Poincaré duality pairs the eigenvalues on \(H^i_\ell\) with those on \(H^{2d-i}_\ell\) by \(\alpha\mapsto q^d/\alpha\). Substituting \(T\mapsto1/(q^dT)\) in the product formula therefore exchanges the corresponding factors up to a monomial, and Corollary \bpout{12.33} carries out this computation and gives (4) in \(\mathbf Q_\ell(T)\). Since the zeta function is rational over \(\mathbf Q\), and the scalar \(\varepsilon q^{d\chi/2}\) is a rational number, the identity holds in \(\mathbf Q(T)\).
\end{proof}

% text: prose-19-09
\phantomsection\label{bp26}%
In the following examples we read the theorem backwards. For projective space the point count determines \(Z\), and the grouping of the reciprocal roots by absolute value recovers the \(P_i\). For an elliptic curve the theorem turns one point count into an inequality.

% object: example-zeta-projective-space-elliptic-curve (blueprint/chap-Proof-of-the-Weil-conjectures/sec-Integrality-and-independence-of-ell/example-zeta-projective-space-elliptic-curve/_section.md)
\phantomsection\label{bp27}%
\begin{bpstatement}{Example~19.8}{ (Projective space and elliptic curves)}
We have \(\#\mathbf P^d(\mathbf F_{q^m})=1+q^m+\cdots+q^{dm}\), so that \(Z(\mathbf P^d_{\mathbf F_q},T)=\prod_{i=0}^d(1-q^iT)^{-1}\). Grouping the reciprocal roots by absolute value as in (2), we find \(P_{2i}=1-q^iT\) for \(0\le i\le d\) and \(P_i=1\) for odd \(i\). For an elliptic curve \(C_0\), we have \(P_0=1-T\) and \(P_2=1-qT\) by (3). Since \(b_1=2\) (Appendix \bpout{D}), the polynomial \(P_1\) has two reciprocal roots \(\alpha,\beta\). Poincaré duality matches them as \(\beta=q/\alpha\), so \(P_1=1-aT+qT^2\) with \(a=\alpha+\beta\). The point-counting formula for \(m=1\) gives \(\#C_0(\mathbf F_q)=1-a+q\), that is, \(a=1+q-\#C_0(\mathbf F_q)\in\mathbf Z\). By (2) both \(\alpha\) and \(\beta\) have absolute value \(q^{1/2}\), and we get Hasse's theorem \(|a|\le2q^{1/2}\).
\end{bpstatement}

% object: example-running-example-purity (blueprint/chap-Proof-of-the-Weil-conjectures/sec-Integrality-and-independence-of-ell/example-running-example-purity/_section.md)
\phantomsection\label{bp28}%
\begin{bpstatement}{Example~19.9}{ (The running example: purity)}
For the curve \(C_0\colon y^2=x^3-x\) over \(\mathbf F_3\) of the introduction, the reciprocal roots \(\pm i\sqrt3\) of \(P_1=1+3T^2\) (Example \bpout{12.32}) have absolute value \(\sqrt3\), as (2) predicts. What purity adds to the earlier examples is a bound: for every genus-one curve over \(\mathbf F_3\) we get \(|4-N_1|\le2\sqrt3\), that is, \(1\le N_1\le7\). In the proof, the upper bound \(|\alpha|\le\sqrt3\) comes from the weight bound for \(Rf_!\) (Theorem \bpout{13.11}) applied to the constant sheaf, and \(\alpha\beta=3\) (Example \bpout{10.23}) turns it into the lower bound.
\end{bpstatement}

% text: prose-19-10
\phantomsection\label{bp29}%
The argument of Hasse's theorem works in every dimension. The point-counting formula expresses \(\#X_0(\mathbf F_{q^m})\) through the traces of \(F^m\), and purity bounds each trace by the dimension of the group times the common absolute value of its eigenvalues.

% object: corollary-point-count-estimates (blueprint/chap-Proof-of-the-Weil-conjectures/sec-Integrality-and-independence-of-ell/corollary-point-count-estimates/_section.md)
\phantomsection\label{bp30}%
\begin{bpstatementit}{Corollary~19.10}{ (Point-count estimates)}
Let \(X_0/\mathbf F_q\) be smooth, proper and geometrically connected of dimension \(d\geq1\), and let \(b_i\) be as in Theorem \hyperref[bp24]{19.7}. Then for every \(m\geq1\),

\[
\left|\#X_0(\mathbf F_{q^m})-(1+q^{dm})\right|
\leq\sum_{i=1}^{2d-1}b_iq^{mi/2}.
\]

For a smooth proper geometrically connected curve of genus \(g\), this is

\[
\left|\#X_0(\mathbf F_{q^m})-(1+q^m)\right|
\leq 2gq^{m/2}.
\]
\end{bpstatementit}

% proof: proof/corollary-point-count-estimates
\begin{proof}[{Proof}]\phantomsection\label{bp31}%
Since \(X_0\) is proper, the point-counting formula of Corollary \bpout{12.31} reads

\[
\#X_0(\mathbf F_{q^m})
=\sum_{i=0}^{2d}(-1)^i\operatorname{Tr}(F^m\mid H^i_\ell).
\]

The terms in degrees \(0\) and \(2d\) are \(1\) and \(q^{dm}\) by Theorem theorem-weil-conjectures-smooth-proper-schemes. The left side of the estimate is therefore the absolute value of the sum of the remaining terms. In degree \(i\), the trace is a sum of \(b_i\) eigenvalues of \(F^m\), each of complex modulus \(q^{mi/2}\) by Theorem theorem-weil-conjectures-smooth-proper-schemes. Its absolute value is thus at most \(b_iq^{mi/2}\), and the triangle inequality proves the first bound. For a curve, \(b_1=2g\) by the cohomology calculation of Appendix \bpout{D}.
\end{proof}



\clearpage\phantomsection\hypertarget{bp-outside}{}%
\section*{Outside this export}
\addcontentsline{toc}{section}{Outside this export}
References marked \ensuremath{{}^\ast} are to parts of the blueprint this document leaves out:
\begin{description}
\item[Corollary 10.13] Degreewise duality over fields, in Ⅱ Duality and the trace formula (\bpslug{corollary-{}\allowbreak{}degreewise-{}\allowbreak{}duality-{}\allowbreak{}over-{}\allowbreak{}fields})
\item[Theorem 10.15] The trace isomorphism in top degree, in Ⅱ Duality and the trace formula (\bpslug{theorem-{}\allowbreak{}trace-{}\allowbreak{}isomorphism-{}\allowbreak{}top-{}\allowbreak{}degree})
\item[Theorem 10.21] Poincaré duality and Frobenius, in Ⅱ Duality and the trace formula (\bpslug{theorem-{}\allowbreak{}poincare-{}\allowbreak{}duality-{}\allowbreak{}frobenius})
\item[Example 10.23] The running example: duality, in Ⅱ Duality and the trace formula (\bpslug{example-{}\allowbreak{}running-{}\allowbreak{}example-{}\allowbreak{}duality})
\item[Chapter 12] The trace formula, in Ⅱ Duality and the trace formula (\bpslug{chap-{}\allowbreak{}The-{}\allowbreak{}trace-{}\allowbreak{}formula})
\item[Corollary 12.31] Point counting and the zeta function, in Ⅱ Duality and the trace formula (\bpslug{corollary-{}\allowbreak{}point-{}\allowbreak{}counting-{}\allowbreak{}zeta-{}\allowbreak{}function})
\item[Example 12.32] The running example: the trace formula, in Ⅱ Duality and the trace formula (\bpslug{example-{}\allowbreak{}running-{}\allowbreak{}example-{}\allowbreak{}trace-{}\allowbreak{}formula})
\item[Corollary 12.33] Functional equation, in Ⅱ Duality and the trace formula (\bpslug{corollary-{}\allowbreak{}functional-{}\allowbreak{}equation})
\item[Theorem 13.11] Weights of $Rf_!$, in Ⅲ Weights and the Riemann hypothesis (\bpslug{theorem-{}\allowbreak{}weights-{}\allowbreak{}rf-{}\allowbreak{}shriek})
\item[Theorem 18.15] Direct image of a mixed sheaf, in Ⅲ Weights and the Riemann hypothesis (\bpslug{theorem-{}\allowbreak{}direct-{}\allowbreak{}image-{}\allowbreak{}mixed-{}\allowbreak{}sheaf})
\item[Appendix D] Cohomology of torsion sheaves on curves, in Appendices (\bpslug{chap-{}\allowbreak{}Cohomology-{}\allowbreak{}of-{}\allowbreak{}torsion-{}\allowbreak{}sheaves-{}\allowbreak{}on-{}\allowbreak{}curves})
\end{description}
\clearpage\phantomsection%
\section*{References}
\addcontentsline{toc}{section}{References}
\bpbibitem{BBD82}{bp32}{A. Beilinson, J. Bernstein, and P. Deligne. \emph{Faisceaux pervers}, volume 100 of \emph{Astérisque}. 1982.}

\bpbibitem{BS15}{bp33}{B. Bhatt and P. Scholze. The pro-étale topology for schemes. \emph{Astérisque}, 369:99–201, 2015. \href{https://people.mpim-bonn.mpg.de/scholze/proetale.pdf}{Author's manuscript}.}

\bpbibitem{Con07}{bp34}{B. Conrad. Deligne's notes on Nagata compactifications. \emph{Journal of the Ramanujan Mathematical Society}, 22:205–257, 2007. \href{https://math.stanford.edu/\~conrad/papers/nagatafinal.pdf}{Author's manuscript}.}

\bpbibitem{Del74}{bp35}{P. Deligne. La conjecture de Weil. I. \emph{Publ. Math. IHÉS}, 43:273–307, 1974.}

\bpbibitem{Del80}{bp36}{P. Deligne. La conjecture de Weil. II. \emph{Publ. Math. IHÉS}, 52:137–252, 1980.}

\bpbibitem{Den22}{bp37}{C. Deninger. There is no “Weil-”cohomology theory with real coefficients for arithmetic curves. \emph{Ann. Sc. Norm. Super. Pisa Cl. Sci.}, 2022. To appear. arXiv:\href{https://arxiv.org/abs/2204.02714}{2204.02714}.}

\bpbibitem{Deu41}{bp38}{M. Deuring. Die Typen der Multiplikatorenringe elliptischer Funktionenkörper. \emph{Abh. Math. Sem. Univ. Hamburg}, 14:197–272, 1941.}

\bpbibitem{Eke90}{bp39}{T. Ekedahl. On the adic formalism. In \emph{The Grothendieck Festschrift, Vol. II}, volume 87 of \emph{Progress in Mathematics}, pages 197–218. Birkhäuser, 1990.}

\bpbibitem{Has36}{bp40}{H. Hasse. Zur Theorie der abstrakten elliptischen Funktionenkörper. I, II, III. \emph{J. Reine Angew. Math.}, 175:55–62, 69–88, 193–208, 1936.}

\bpbibitem{Hum72}{bp41}{J. E. Humphreys. \emph{Introduction to Lie Algebras and Representation Theory}, volume 9 of \emph{Graduate Texts in Mathematics}. Springer, 1972.}

\bpbibitem{Ill94}{bp42}{L. Illusie. Autour du théorème de monodromie locale. \emph{Astérisque}, 223:9–57, 1994. §4. \href{https://www.numdam.org/item/AST_1994__223__9_0/}{Original article}.}

\bpbibitem{Kat88}{bp43}{N. M. Katz. Travaux de Laumon. \emph{Astérisque}, 161–162:105–132, 1988. Séminaire Bourbaki, exp. 691.}

\bpbibitem{KW01}{bp44}{R. Kiehl and R. Weissauer. \emph{Weil conjectures, perverse sheaves and \(\ell\)-adic Fourier transform}, volume 42 of \emph{Ergebnisse der Mathematik}. Springer, 2001.}

\bpbibitem{Lau87}{bp45}{G. Laumon. Transformation de Fourier, constantes d'équations fonctionnelles et conjecture de Weil. \emph{Publ. Math. IHÉS}, 65:131–210, 1987.}

\bpbibitem{Mil86a}{bp46}{J. S. Milne. \href{https://www.jmilne.org/math/xnotes/AVs.pdf}{Abelian Varieties}. In \emph{Arithmetic Geometry}. 1986. Corrected online version.}

\bpbibitem{Mil86b}{bp47}{J. S. Milne. \href{https://www.jmilne.org/math/xnotes/JVs.pdf}{Jacobian Varieties}. In \emph{Arithmetic Geometry}. 1986. Corrected online version.}

\bpbibitem{Ran39}{bp48}{R. A. Rankin. Contributions to the theory of Ramanujan's function \(\tau(n)\) and similar arithmetical functions. II. \emph{Proc. Cambridge Philos. Soc.}, 35:357–372, 1939.}

\bpbibitem{SGA4}{bp49}{M. Artin, A. Grothendieck, and J.-L. Verdier, editors. \emph{Théorie des topos et cohomologie étale des schémas}, volumes 269, 270, 305 of \emph{Lecture Notes in Mathematics}. Springer, 1972–1973. SGA 4.}

\bpbibitem{SGA4½}{bp50}{P. Deligne. \emph{Cohomologie étale}, volume 569 of \emph{Lecture Notes in Mathematics}. Springer, 1977. SGA 4½; \emph{Théorèmes de finitude}, §3, and \emph{Rapport sur la formule des traces}, §3, and \emph{Sommes trigonométriques}, §1.}

\bpbibitem{SGA5}{bp51}{A. Grothendieck, editor. \emph{Cohomologie \(\ell\)-adique et fonctions \(L\)}, volume 589 of \emph{Lecture Notes in Mathematics}. Springer, 1977. SGA 5.}

\bpbibitem{SGA7II}{bp52}{P. Deligne and N. Katz, editors. \emph{Groupes de monodromie en géométrie algébrique}, volume 340 of \emph{Lecture Notes in Mathematics}. Springer, 1973. SGA 7 II; Exposé XIII, (2.1.7).}

\bpbibitem{Stacks}{bp53}{The Stacks Project Authors. \href{https://stacks.math.columbia.edu}{The Stacks Project}. Individual results are cited by their stable tags in the text.}

\bpbibitem{Var05}{bp54}{Y. Varshavsky. Lefschetz–Verdier trace formula and a generalization of a theorem of Fujiwara. 2005. Version 2. arXiv:\href{https://arxiv.org/abs/math/0505564}{math/0505564}. \href{https://arxiv.org/abs/math/0505564v2}{Author's preprint}.}

\bpbibitem{Wei48a}{bp55}{A. Weil. \emph{Sur les courbes algébriques et les variétés qui s'en déduisent}. Hermann, Paris, 1948.}

\bpbibitem{Wei48b}{bp56}{A. Weil. \emph{Variétés abéliennes et courbes algébriques}. Hermann, Paris, 1948.}

\bpbibitem{Wei49}{bp57}{A. Weil. Numbers of solutions of equations in finite fields. \emph{Bull. Amer. Math. Soc.}, 55:497–508, 1949.}

\end{document}
